\section{Discussion}

\label{sec:discussion}
% ══════════════════════════════════════════════════════════════════════════════

\subsection{Why Adaptive Gamification Works}

The key insight is that gamification should optimize for scientific value,
not just engagement. Standard gamification (leaderboards, streaks)
increases observation volume but not necessarily data value. Adaptive
gamification using RL-generated challenges aligns individual reward
structures with collective scientific needs.

The reinforcement learning formulation is crucial: it balances the
exploration-exploitation tradeoff between directing users to known data
gaps (exploitation) and discovering new user preferences and capabilities
(exploration). Thompson sampling naturally manages this tradeoff through
posterior uncertainty.

\subsection{The Role of AI Assistance}

AI-assisted identification increases accuracy from 78.4\% to 91.3\%,
but the anchoring mitigation strategies are essential. Without delayed
display and disagreement incentives, AI assistance reduces identification
diversity---users simply accept the AI suggestion without independent
judgment. Our design preserves the ``human in the loop'' while leveraging
AI to support novice observers.

\subsection{Token Incentives vs.\ Intrinsic Motivation}

A central concern with token incentives is that extrinsic rewards may
crowd out intrinsic motivation \citep{deci1999meta}. Our user surveys
indicate that token incentives primarily attract new users (42\% cite
tokens as initial motivation), while retained users are primarily
motivated by learning (67\%), contribution to science (54\%), and
community belonging (48\%). Tokens serve as an on-ramp to intrinsic
engagement rather than a sustained motivator.

\subsection{Limitations}

\begin{enumerate}
  \item \textbf{Smartphone dependency}: ZooPlatform requires a
    smartphone with GPS and camera, excluding potential contributors
    in low-connectivity regions.
  \item \textbf{Species coverage}: Despite taxonomic bias reduction,
    invertebrates and fungi remain underrepresented due to
    identification difficulty.
  \item \textbf{Night observations}: The platform underperforms for
    nocturnal species despite Night Watch challenges, as photo quality
    in low light is poor.
  \item \textbf{Token speculation}: If ZooRep became tradeable (it is
    currently non-transferable), gaming incentives would increase
    substantially.
  \item \textbf{Privacy}: GPS-tagged observations may reveal observer
    location patterns. We implement location fuzzing for human activity
    but not for biodiversity observations.
\end{enumerate}


% ══════════════════════════════════════════════════════════════════════════════
